% Transcription — fonds Alexandre Grothendieck, shelfmark 143, batch 1 (pages 1-20).
% Established from the facsimile archives/batches/143/batch-01.pdf, cut from a
% copy of 143.pdf supplied by hand (the archive host could not be reached from
% the session); PDF page k is the archivists' page k, no cover sheet. Per the
% skill transcribe-grothendieck. The folder is thirty-three pages; this is the
% first of two batches (pages 21-33 are batch 2, transcribed in parallel by
% another agent and read at the end of this pass for notation).
%
% Pass: Opus 5.5 (claude-opus-5-5), 2026-09-24 — first pass, unchecked against the pages by a human.
%
% Pages skipped: none. Pages summarised: none (no typescript). All twenty
% pages are his, pencil on pages 1-4, blue ink from page 4 on (the ink
% changes mid-page 4), on small white sheets; page 19 and the top of page 20
% are crossed by two long parallel oblique strokes, apparently a cancellation,
% and are transcribed under a \note.
%
% His own pagination: 1 (p. 1), 2 (p. 3), 3 (p. 5), 4 (p. 7), 5 (p. 9),
% 6 (p. 11), 7 (p. 13), 8 (p. 15), 9 (p. 17), 10 (p. 19) — every recto; the
% even archivists' pages are the versos and carry no number. This runs on into
% batch 2 (p. 21 = his 11). No title of his on page 1; the section heading is
% the edition's and says so. He numbers his paragraphs by circled numerals
% 1-8 (kept as \textcircled) and a recapitulation by roman I, II, III; his
% formula numbers restart at (1) on page 9.
%
% What the batch is about. The group-theoretic set-up of the « paradigmes »:
% for a scheme X-bar of finite type over Q-bar (then a script X), the group
% G_X of automorphisms of X as an absolute scheme, 1 -> G -> G_X -> Gamma =
% Gal(Q-bar/Q) with image of finite index; a universal cover X~ and
% Sigma_X = Aut(X~ // X), 1 -> pi_1(X) -> Sigma_X -> G_X -> 1 (6); a point xi
% in X(Q-bar), its stabiliser G_X(xi), a partial splitting (8) and a germ of
% lifting phi_xi : Gamma' -> Sigma_X on an open subgroup (9)-(11), the
% remark that xi is determined by its germ when X embeds in a semi-abelian
% variety (§ 3). § 4: a reduced effective divisor D = sum D_i, X' = X - D,
% the subgroup G_X' preserving D, the extension (14), G_X' -> S_I (16), the
% Kummer homomorphisms phi_i : T(Q-bar) = lim mu_n -> pi_1(X') (17), the
% cyclotomic character chi (18), whence G_X' -> Autext(pi_1(X'), T, (phi_i))
% (19); the normalisation of X in X~', the sets X~'(I), X~'(i) and the maps
% Phi_i, Phi : X~'(I) -> Hom(T(Q-bar), pi_1(X')) (22)-(23). § 5: a set J of
% rational points, the group G' preserving D and J, X~'(J) -> J (24) and the
% partial splittings Psi (25). § 6 « Récapitulons » (p. 8-10): data I
% (X, D, J), derived II (G', G', T, Gamma), choice III (a universal cover),
% and the structures (1)-(9) — extension, Sigma'-sets X~'(I), X~'(J),
% Phi, the system Psi of splittings, the germ Gamma' -> Sigma'. Page 11: the
% « Programme des Gleichnis »: (a) a complete group-theoretic description of
% Gamma; (b) a complete description of all the structure above in typical
% cases, and an NB on passing to finite étale covers (« bouclage » /
% « fixage de tresses »). Page 12-13: for Gamma' fixed by a descent to a
% number field K, the extension Sigma'_Gamma' (9), bitorsors B_gamma acting on
% Revét(X') = Ensf(pi_1) by E |-> B_gamma ∧ E (10), and the resulting
% isomorphisms (11) read as the action of Sigma' on X~'(I), X~'(J). § 7
% (p. 14-16): reconstruction of the fundamental groupoid Pi_1(X'; J) from
% X~'(J) with its free right pi_1-action, Isom_Pi1(j, j') = Isom of
% pi_1-torsors, and of the action of G' through the bitorsors B_g; then a
% « Reformulation »: an equivalence alpha / beta between the category C of
% (1 -> pi -> Sigma -> G -> 1, Z a Sigma-set with pi acting freely) and the
% category C' of (Pi, Pi~, p, G, action), defined on p. 17-18. § 8 (p. 18-20):
% a triple (T, T, Phi) — a Sigma-set script T, a group T (≃ Z^ ?) with chi,
% Phi : T -> Hom_gr(T, pi) — with J = T/pi and a groupoid <T, J> mapped to
% the groupoid Pi^ of universal covers (struck through); then the groupoid
% <T, pi> with pi_0 = T/pi = J, pi_1(T, t) = pi_t, T commutative and
% Phi_t : T ~ pi_t, breaking off mid-sentence into batch 2.
%
% Notation fixed once for the batch. Script X as \mathfrak{X} (as he writes
% it, a round X; X-bar on p. 1 before he switches); \mathcal{G} for his
% script G (G_X, G', G_j), plain G for G_X-bar = Aut(X-bar/Q-bar) and G'
% (p. 1, 8, 9); from the « Reformulation » (p. 16) the same round G is set
% plain G, as batch 2 does. The extension letter, which wavers on the page
% between a round E and a sigma, is set \Sigma throughout, as in batch 2.
% \mathfrak{Z} for the curly « 3 »-shaped letter of pp. 16-18 (the Sigma-set
% corresponding to X~'(J)); \mathcal{T} for the barred round T of pp. 18-20
% (corresponding to X~'(I)), as in batch 2, and roman T for T(Q-bar) and the
% group T. \mathcal{B}_gamma, \mathcal{B}_g for his round B (bitorsors);
% \mathfrak{S}_I; \wedge^{\pi_1} for his contracted product ∧ with pi_1
% under the wedge.
%
% Batch-2 cross-check. Batch 2's \Sigma is the same letter as this batch's
% extension letter (set \Sigma here to match). Batch 2's 𝒯 and <𝒯, π> agree
% with pp. 18-20 here; page 20 opens <𝒯, π> exactly as page 21 continues it.
%
% Hardest pages: 1 (pencil; the oblique margin notes), 4 (the germ formula
% (11) and its blacked-out argument), 6 and 7 (the long oblique margin note
% of p. 7 is almost wholly illegible), 10 (the boxed margin note and the
% struck (8)), 11 (the German words and the NB), 17 (layout). Readings to
% check first: « groupe de base » (p. 1); « $/\!/$ » in (5) (p. 1); the whole
% of (11) (p. 4); « circonstanciel » (p. 5); the second argument of Autext in
% (19) (p. 6); « gruppentheoretisch », « bouclage de tresses », « fixage de
% tresses » (p. 11); « RemB » (p. 12); « Ob Pi_1 = J » (p. 14); the words
% under the two oblique strokes (p. 19); the phrase before « T commutatif »
% (p. 20).
\documentclass[11pt,a4paper]{article}
% Le préambule commun à toutes les transcriptions du fonds Grothendieck.
%
% Il tient en une idée : **l'incertitude se compose, elle ne se résout pas.**
% Une transcription qui lisse un mot douteux a détruit la seule chose pour
% laquelle on la faisait — un lecteur doit pouvoir savoir, à chaque endroit,
% ce qui était sur le feuillet et ce qui a été ajouté. Les six macros qui
% suivent sont donc l'appareil critique, et non de la décoration.
%
% Les mêmes commandes sont reconnues par `scripts/render.mjs`, qui produit la
% vue de lecture du site. Une macro ajoutée ici doit l'être aussi là-bas,
% faute de quoi le rendu échouera bruyamment — ce qui est voulu : un rendu
% silencieusement faux est pire qu'un rendu absent.

\ProvidesPackage{grothendieck}[2026/08/08 Transcriptions du fonds Grothendieck]

\usepackage{amsmath,amssymb,amsthm}
\usepackage{mathrsfs}
\usepackage{stmaryrd}
% Les diagrammes commutatifs sont partout dans ce fonds : c'est la langue
% même de la géométrie algébrique de Grothendieck, et une transcription qui
% les rendrait en prose ne transcrirait rien.
\usepackage{tikz-cd}
% Français par défaut — c'est la langue des feuillets — mais l'anglais est
% chargé aussi : la traduction et le résumé sont des documents anglais, et la
% typographie française (espaces avant les deux-points) y serait une faute.
% Ces documents basculent par \selectlanguage{english} en tête de corps.
\usepackage[english,french]{babel}
\usepackage{fontspec}
\usepackage{geometry}
\geometry{a4paper,margin=25mm,marginparwidth=28mm}
\usepackage{marginnote}
\usepackage{xcolor}
% ulem, not soul. `soul`'s \st and \ul are built on a character-by-character
% scan that XeLaTeX's font machinery breaks: the document fails with an
% undefined control sequence rather than degrading. ulem does the same two jobs
% and is Unicode-engine safe. `normalem` stops it hijacking \emph, which the
% transcriptions use for Grothendieck's own emphasis.
\usepackage[normalem]{ulem}
\usepackage{hyperref}
\hypersetup{colorlinks=true,linkcolor=black,urlcolor=black,pdfborder={0 0 0}}

% --- Métadonnées du lot -----------------------------------------------------
% Reprises telles quelles par le script de rendu, qui les affiche en tête de
% la vue de lecture. Elles fixent ce que le fichier prétend couvrir : sans
% elles, une transcription flotte sans point d'attache dans un fonds de seize
% mille feuillets.

\newcommand{\folder}[1]{\gdef\grothFolder{#1}}
\newcommand{\batch}[1]{\gdef\grothBatch{#1}}
\newcommand{\leaves}[2]{\gdef\grothFirst{#1}\gdef\grothLast{#2}}
\newcommand{\foldertitle}[1]{\gdef\grothFoldertitle{#1}}
\gdef\grothFolder{?}\gdef\grothBatch{?}\gdef\grothFirst{?}\gdef\grothLast{?}\gdef\grothFoldertitle{}

% --- Le résumé --------------------------------------------------------------
%
% Il ouvre la lecture modernisée, sous le titre « Résumé », et il est composé
% en petit. Ce n'est pas un ornement : c'est le seul passage du document qui
% ne soit pas des mathématiques, et un lecteur doit pouvoir le reconnaître
% avant de l'avoir lu — puis le sauter s'il connaît déjà le sujet, ou s'y
% arrêter s'il ne le connaît pas. Le filet qui le suit marque la reprise du
% corps.

\newenvironment{resume}
  {\small\setlength{\parskip}{0.45em}}
  {\par\medskip\hrule\medskip}

% --- Mots-clés ---------------------------------------------------------------
%
% En anglais, en clôture du résumé de la lecture modernisée : le vocabulaire
% moderne sous lequel chercher ce que les feuillets construisent. Le manifeste
% du site (`npm run manifest`) les extrait pour étiqueter la cote entière —
% cette ligne est la source unique des tags, il n'y a pas de fichier à côté.
\newcommand{\keywords}[1]{\par\smallskip\noindent
  {\footnotesize\textsc{Keywords} --- \textit{#1}}\par}

% --- L'appareil critique ----------------------------------------------------

% Le numéro de feuillet, en marge. C'est l'ancre qui permet de régler un
% désaccord sur une formule en regardant *un* feuillet plutôt qu'une chemise
% de six cents. Le site s'en sert aussi pour tourner les pages du fac-similé
% quand on fait défiler la transcription.
\newcommand{\leaf}[1]{%
  \marginnote{\footnotesize\color{gray!70}\textsf{#1}}%
  \label{leaf:#1}\ignorespaces}

% Illisible : on ne devine pas. Un mot inventé qui se lit comme les autres est
% le pire résultat possible.
\newcommand{\ill}{\textcolor{red!60!black}{[\,\ldots\,]}}

% Lecture douteuse : le mot est donné, et signalé comme incertain.
\newcommand{\uncertain}[1]{\uline{#1}}

% Ajout de l'éditeur — une lettre manquante, un mot sous-entendu.
\newcommand{\add}[1]{\textcolor{blue!50!black}{[#1]}}

% Biffé par Grothendieck lui-même. Ce qu'il a rayé fait partie du document :
% une rature dit souvent où la pensée a changé de direction.
\newcommand{\struck}[1]{\sout{#1}}

% Note du transcripteur, sur l'état matériel du feuillet.
\newcommand{\note}[1]{\par\smallskip{\small\color{gray!60!black}\textit{[#1]}}\par\smallskip}

% Note marginale de Grothendieck, distincte du corps du texte.
\newcommand{\marginal}[1]{\par\smallskip\noindent\hspace{1em}%
  {\small\color{gray!60!black}#1}\par\smallskip}

% --- Titre ------------------------------------------------------------------

\AtBeginDocument{%
  \begin{center}
    {\large\bfseries Fonds Alexandre Grothendieck --- cote n\textsuperscript{o}~\grothFolder}\\[2pt]
    {\itshape\grothFoldertitle}\\[4pt]
    {\small feuillets \grothFirst--\grothLast\ (lot \grothBatch)}\\[2pt]
    {\footnotesize\color{gray!60!black}Transcription établie d'après le fac-similé numérisé
     par l'Université de Montpellier.\\
     Les lectures incertaines sont soulignées, les illisibles notées [\,\ldots\,],
     les ajouts entre crochets.}
  \end{center}
  \vspace{1em}}

\endinput


\folder{143}
\batch{1}
\pages{1}{20}
\dating{[vers 1980-1981]}
\watermark{Édition de démonstration}
\foldertitle{Paradigmes des courbes algébriques (version provisoire) : notes manuscrites (s.d.)}

\begin{document}

\section{Le groupe $\mathcal{G}_{\overline{X}}$ et l'extension $\Sigma_{\mathfrak{X}}$}

\note{ce titre est de l'édition ; la page 1 ne porte pas de titre de sa main.
Il numérote ses paragraphes par des chiffres cerclés et ses formules par
(1), (2), \ldots}

\page{1}
\note{p.~1 de l'auteur}
\textcircled{1} Soit $\overline{\mathbb{Q}}$ la clôture algébrique de $\mathbb{Q}$ dans
$\mathbb{C}$. Si \struck{$X$} $\overline{X}$ est un schéma sur
$\overline{\mathbb{Q}}$, on considère le groupe $\mathcal{G}_{\overline{X}}$
des automorphismes de $\overline{X}$ en tant que schéma absolu. On a une suite
\[
  (1)\qquad 1 \to G_{\overline{X}} \longrightarrow \mathcal{G}_{\overline{X}}
  \longrightarrow \Gamma
\]
\[
  (2)\qquad \left\{
  \begin{array}{l}
    \Gamma = \mathrm{Gal}(\overline{\mathbb{Q}}/\mathbb{Q}) \\
    G_{\overline{X}} = \mathrm{Aut}(\overline{X}/\overline{\mathbb{Q}})
  \end{array}\right.
\]
\marginal{en marge gauche, au crayon, en oblique, à hauteur de (1) :
« Question : \ill{} \uncertain{d'indice fini} ? \ill{} \uncertain{vrai} »}
\marginal{en marge droite, au crayon, en oblique :
$\overline{\mathbb{Q}} \hookrightarrow \Gamma(\mathfrak{X}, \mathcal{O}_{\mathfrak{X}})$,
$\mathbb{Q} \subset \overline{\mathbb{Q}}$ ; « il faut \ill{} p.~ex. \ill{}
$\overline{\mathbb{Q}}$ \ill{} de $\Gamma_{\mathfrak{X}}$ »}

(3) Les sous-extensions $K$ de $\overline{\mathbb{Q}}/\mathbb{Q}$ correspondent
aux sous-groupes fermés $\Gamma_K$ de $\Gamma$ (d'indice fini ssi $K$ fini sur
$\mathbb{Q}$). À une restriction du \uncertain{groupe} de base
$\overline{\mathbb{Q}}$ de $\overline{X}$ à $K$ correspond un relèvement
\note{« groupe de base » est lu tel quel ; on attendrait « corps de base »}

\begin{tikzcd}[column sep=large]
(4)\quad \mathcal{G}_{\overline{X}} \arrow[r] & \Gamma \\
 & \Gamma_K \arrow[ul] \arrow[u, hook]
\end{tikzcd}

qui le définit si $\overline{X}$ est réduit \marginal{un « ? » cerclé, au
crayon, sur « réduit »} et loc.\ de t.f.\ sur $\overline{\mathbb{Q}}$. Cela
montre que si $\overline{X}$ est de t.f.\ \struck{\ill{}} sur
$\overline{\mathbb{Q}}$, l'hom.\ $\mathcal{G}_{\overline{X}} \to \Gamma$ a une
image d'indice fini (puisqu'on \uncertain{peut toujours} réduire le corps de
base à une ext.\ finie $K$ de $\mathbb{Q}$).
\marginal{en marge gauche, au crayon, en oblique, près de (4) : « grain de sel
d'effectivité \ill{} de \uncertain{projectivité} \ill{} »}

\textcircled{2} Supposons $\mathfrak{X}$ connexe $\neq \emptyset$, \add{et}
\struck{considérons un} revêtement universel
$\widetilde{\mathfrak{X}} \to \mathfrak{X}$ de $\mathfrak{X}$, et
\struck{considérons} posons
\[
  (5)\qquad \Sigma(\mathfrak{X}, \widetilde{\mathfrak{X}})
  \ \text{ou}\ \Sigma_{\mathfrak{X}}
  = \mathrm{Aut}(\widetilde{\mathfrak{X}} \,/\!/\, \mathfrak{X})
\]
\note{le signe entre $\widetilde{\mathfrak{X}}$ et $\mathfrak{X}$ est une barre
oblique surchargée, rendue « $/\!/$ » sans préjuger de sa lecture.
Le passage de $\overline{X}$ à $\mathfrak{X}$ (ronde) est de lui. La lettre
rendue $\Sigma$, ici et dans tout le lot, hésite sur la page entre un E ronde
et un sigma ; elle est rendue $\Sigma$ comme dans le second lot (p.~21 sqq.)}

\page{2}
On a donc une suite exacte
\[
  (6)\qquad 1 \longrightarrow \pi_1(\mathfrak{X}) \longrightarrow
  \Sigma_{\mathfrak{X}} \longrightarrow \mathcal{G}_{\mathfrak{X}}
  \longrightarrow 1 ,
  \qquad \pi_1(\mathfrak{X}) = \mathrm{Aut}(\widetilde{\mathfrak{X}}/\mathfrak{X}),
  \qquad \mathcal{G}_{\mathfrak{X}} \to \Gamma
\]
\note{dans l'original, « $=\mathrm{Aut}(\widetilde{\mathfrak{X}}/\mathfrak{X})$ »
est écrit sous $\pi_1(\mathfrak{X})$ avec un signe $=$ vertical, et une flèche
descend de $\mathcal{G}_{\mathfrak{X}}$ vers $\Gamma$}

\struck{La donnée d'un point} Considérons un point $\xi$
\[
  (7)\qquad \xi \in \mathfrak{X}(\overline{\mathbb{Q}})
\]
et soit $\mathcal{G}_{\mathfrak{X}}(\xi) \subset \mathcal{G}_{\mathfrak{X}}$ le
sous-groupe de $\mathcal{G}_{\mathfrak{X}}$ qui fixe $\xi$. (NB. Si
\marginal{interligne, souligné : « $\mathfrak{X}$ est de t.f.\ $/\overline{\mathbb{Q}}$ et »}
$G_{\mathfrak{X}}$ est fini, $\mathcal{G}_{\mathfrak{X}}$ est profini et
$\mathcal{G}_{\mathfrak{X}}(\xi)$ est un s-groupe ouvert\ldots). On trouve alors,
dès qu'on choisit pour $\widetilde{\mathfrak{X}}$ le rev.\ universel
\marginal{interligne : $\widetilde{\mathfrak{X}}(\xi)$} pointé au dessus de
$\xi$, un splitting partiel de (6)

\begin{tikzcd}
(8)\quad \Sigma_{\mathfrak{X}} \arrow[r] & \mathcal{G}_{\mathfrak{X}} \arrow[r] & 1 \\
 & \mathcal{G}_{\mathfrak{X}}(\xi) \arrow[ul] \arrow[u, hook]
\end{tikzcd}

\note{le trait vertical entre $\mathcal{G}_{\mathfrak{X}}$ et
$\mathcal{G}_{\mathfrak{X}}(\xi)$ est rendu comme l'inclusion ; le sens du trait
n'est pas net sur la page}

Dans le cas où $\widetilde{\mathfrak{X}}$ est quelconque, on trouve \uncertain{qu'un}
hom.\ mod.\ composition avec \uncertain{automorphisme intérieur} provenant de
$\pi_1(\mathfrak{X})$. \struck{Si $\mathfrak{X}$ est de t.f.\ et}

\page{3}
\note{p.~2 de l'auteur}
\struck{s'envoie injectivement dans un groupe algébrique commutatif de rang
unipotent nul (ext.\ \uncertain{cat.}\ d'une VA par un tore) sous
l'application $\mathfrak{X}(\overline{\mathbb{Q}}) \to J$}
\note{ces lignes, qui continuent la phrase biffée au bas de la page~2, sont
barrées de traits obliques}

\textcircled{3} Supposons $\mathfrak{X}$ de t.f.\ sur $\overline{\mathbb{Q}}$.
Alors on a un « germe de relèvement » de $\Gamma$

\begin{tikzcd}
(9)\quad \mathcal{G}_{\overline{X}} \arrow[r] & \Gamma \\
 & \Gamma' \arrow[ul] \arrow[u, hook]
\end{tikzcd}

sur un s-groupe ouvert (modulo la relation de coïncidence sur un s-groupe
ouvert plus petit). Pour \marginal{interligne : « $\xi \in
\mathfrak{X}(\overline{\mathbb{Q}})$ donné, et pour »} $\Gamma'$ assez petit,
\marginal{interligne : « de $\xi$ »} $\Gamma' \to \mathcal{G}_{\overline{X}}$
se factorise par $\mathcal{G}_{\overline{X}}(\xi)$, donc on trouve un
\marginal{interligne : « germe d' »} hom.\ compatible
\[
  \Gamma' \longrightarrow \Sigma_{\mathfrak{X}}
\]
\struck{$\Gamma' \to \Sigma_{\mathfrak{X}}$}
rendant commutatif
\marginal{en marge gauche, au crayon : $\mathfrak{X} = \overline{X}$}

\begin{tikzcd}
(10)\quad \Sigma_{\mathfrak{X}} \arrow[r] & \mathcal{G}_{\mathfrak{X}} \arrow[r] & \Gamma \\
 & & \Gamma' \arrow[ull, "\varphi_{\xi}"] \arrow[u, hook]
\end{tikzcd}

--- du moins si $\widetilde{\mathfrak{X}} = \widetilde{\mathfrak{X}}(\xi)$ ---
d'ailleurs un hom.\ défini mod.\ composition par un \uncertain{automorphisme
intérieur} provenant de $\pi_1(\mathfrak{X})$. Scholie : \ill{}

\page{4}
encore une fois que c'est le germe seulement de $\varphi_{\xi}$ qui est défini
(mod restriction à s-groupe d'indice fini). [Si $\mathfrak{X}$ s'envoie
injectivement dans une extension d'une VA par un tore, alors $\xi$ est connu
quand on connaît son germe
\[
  (11)\qquad \varphi_{\xi} \in \varinjlim_{\Gamma' \subset \Gamma}
  \mathrm{Hom}(\Gamma', \Sigma_{\mathfrak{X}})
\]
\ill{} $(\pi_1(\mathfrak{X}), \Sigma_{\mathfrak{X}})$
\note{sous la limite, un mot souligné en pointillé, sans doute « ouvert » ;
après $\mathrm{Hom}(\Gamma', \Sigma_{\mathfrak{X}})$ un argument est
noirci et barré, puis vient le groupe écrit ici, dont le signe qui le précède
ne se lit pas}
[et à l'image de $\varphi_{\xi}$ dans un $H^1$ convenable\ldots]]

\textcircled{4}
\[
  (12)\qquad \text{Soit } \boxed{D = \sum_{i \in I} D_i}
\]
\marginal{interligne, au-dessus : « effectif »}
(un diviseur) sur $\mathfrak{X}$, ($I$ l'ens.\ de ses comp.\ irréd.) sans
comp.\ multiples, \marginal{interligne : « \uncertain{i.e.\ lisse} ) »}
supposons que $\mathfrak{X}$ soit régulier \uncertain{en les} pts génériques
de chaque $D_i$. Soit \struck{$\mathfrak{X}_I$}
\[
  (13)\qquad \mathfrak{X}' = \mathfrak{X} - \operatorname{supp} D
\]
et considérons les sous-groupes $\mathcal{G}_{\mathfrak{X};\mathfrak{X}'}$
\marginal{souligné sous l'indice : « \uncertain{ouvert} »} de
$\mathcal{G}_{\mathfrak{X}}$ qui invarie $D$. On les notera (\uncertain{ou},
plus précisément, $\mathcal{G}_{\mathfrak{X}}(I)$, abus de notation)
$\mathcal{G}_{\mathfrak{X}'}$, (\textbf{D} \ill{}

\page{5}
\note{p.~3 de l'auteur}
\[
  (14)\qquad 1 \to \pi_1(\mathfrak{X}') \longrightarrow
  \Sigma_{\mathfrak{X}'} \longrightarrow \mathcal{G}_{\mathfrak{X}'}
  \to 1 .
\]
et on a encore un germe de \struck{hom} \marginal{interligne : « relèvement »}
(\uncertain{circonstanciel} mod comp.\ par autom.\ int.\ provenant de
$\pi_1(\mathfrak{X}')$)

\begin{tikzcd}
(15)\quad \Sigma_{\mathfrak{X}'} \arrow[r] & \mathcal{G}_{\mathfrak{X}'} \arrow[r] & \Gamma \\
 & & \Gamma' \arrow[ull, "\varphi_{\xi}"] \arrow[u, hook]
\end{tikzcd}

(si $\xi \in \mathfrak{X}'(\overline{\mathbb{Q}})$).

Notons qu'on a ici un hom
\[
  (16)\qquad \mathcal{G}_{\mathfrak{X}'} \longrightarrow \mathfrak{S}_I
\]
La théorie de Kummer nous fournit, pour tout $i$, un hom extérieur
\[
  (17)\qquad T(\overline{\mathbb{Q}}) \xrightarrow{\ \varphi_i\ }
  \pi_1(\mathfrak{X}'),
  \qquad T(\overline{\mathbb{Q}}) = \varprojlim \mu_n(\overline{\mathbb{Q}}),
  \qquad \varphi_i \in \mathrm{Homext}(T(\overline{\mathbb{Q}}), \pi_1(\mathfrak{X}'))
\]
\note{la flèche porte une barre verticale en travers ; l'égalité
$T(\overline{\mathbb{Q}}) = \varprojlim \mu_n(\overline{\mathbb{Q}})$ est écrite
sous $T(\overline{\mathbb{Q}})$, et « $\varphi_i \in \mathrm{Homext}(\ldots)$ » à
droite, en marge. À gauche, une note au crayon de plusieurs lignes, où se lit
« (17) », est hachurée de traits verticaux et ne se lit pas}

Notons aussi un hom canonique
\[
  (18)\qquad (\mathcal{G}_{\mathfrak{X}'} \to)\ \Gamma
  \xrightarrow{\ \chi\ } \widehat{\mathbb{Z}}^{*} \xrightarrow{\ \sim\ }
  \mathrm{Aut}(T(\overline{\mathbb{Q}}))
\]
faisant opérer $\mathcal{G}_{\mathfrak{X}'}$ via $\Gamma$ sur
$T(\overline{\mathbb{Q}})$. Ceci dit, l'\struck{opération} opération extérieure
de $\mathcal{G}_{\mathfrak{X}'}$ sur $\pi_1(\mathfrak{X}')$ est compatible avec
la famille

\page{6}
des $(\varphi_i)$ (« structure Kummérienne ») compte tenu des opérations
\marginal{interligne : « (16) et (18) »} de $\mathcal{G}_{\mathfrak{X}'}$ sur
$I$ et $T(\overline{\mathbb{Q}})$. On peut donc écrire
\[
  (19)\qquad \mathcal{G}_{\mathfrak{X}'} \longrightarrow
  \mathrm{Autext}\bigl(\pi_1(\mathfrak{X}'), T(\overline{\mathbb{Q}}), (\varphi_i)\bigr).
\]
\note{le deuxième argument, $T(\overline{\mathbb{Q}})$, est surchargé et de lecture incertaine}

\struck{Supposons $\mathfrak{X}$} Soit $\widehat{\widetilde{\mathfrak{X}}}{}'$
le normalisé de $\mathfrak{X}$ \struck{\ill{}} dans l'« extension » définie par
$\widetilde{\mathfrak{X}}'$, donc on a un morphisme profini
\[
  \widehat{\widetilde{\mathfrak{X}}}{}' \longrightarrow \mathfrak{X}
\]
Soit $\widetilde{\mathfrak{X}}'(I)$ l'image inverse de l'ens.\ des pts
\struck{génériques} \marginal{interligne : « maximaux »} des $D = \sum D_i$.
On a donc
\[
  (20)\qquad \widetilde{\mathfrak{X}}'(I) \longrightarrow I
\]
et une opération de \struck{\ill{}} $\Sigma_{\mathfrak{X}'}$ sur
$\widetilde{\mathfrak{X}}'(I)$ \uncertain{qui transporte la structure)}
compatible avec les opérations sur $I$ et l'application (20). Pour $i \in I$,
soit
\[
  (21)\qquad \widetilde{\mathfrak{X}}'(i) \subset \widetilde{\mathfrak{X}}'(I)
\]
sa fibre dans $\widetilde{\mathfrak{X}}'(I)$, \ill{}
\marginal{à droite de (21) : « (C'est un \uncertain{espace homogène} sous
$\pi_1(\mathfrak{X}')$. On \struck{trouve} \ill{} »}

\page{7}
\note{p.~4 de l'auteur}
une application
\[
  (22)\qquad \widetilde{\mathfrak{X}}'(i) \xrightarrow{\ \Phi_i\ }
  \mathrm{Hom}(T(\overline{\mathbb{Q}}), \pi_1(\mathfrak{X}'))
\]
compatible avec les opérations de $\pi_1(\mathfrak{X}')$ \uncertain{de part} et
d'autre, et qui précise \uncertain{l'hom.\ extérieur} $\varphi_i$ de (17). Mieux
encore, le système $(\Phi_i)_{i \in I} = \Phi$,
\[
  (23)\qquad \widetilde{\mathfrak{X}}'(I) \xrightarrow{\ \Phi\ }
  \mathrm{Hom}(T(\overline{\mathbb{Q}}), \pi_1(\mathfrak{X}'))
\]
est compatible avec les $\Sigma_{\mathfrak{X}'}$ de part et d'autre, en
se rappelant, pour l'opération de $\Sigma_{\mathfrak{X}'}$ sur le but de
droite, de tenir compte aussi de l'opération de $\Sigma_{\mathfrak{X}'}$
sur $T(\overline{\mathbb{Q}})$ (via $\Gamma$).
\marginal{en marge gauche, au crayon, une note de plusieurs lignes obliques,
presque entièrement illisible : « ici il faudrait \ill{} dire \ill{}
$\mathfrak{X}$ \ill{} $\widetilde{\mathfrak{X}}'(I)$ \ill{} $\Phi$ \ill{} »}

\textcircled{5} Soit enfin, par dessus le marché, une partie
$J \subset \mathfrak{X}'(\overline{\mathbb{Q}})$. Soit
$\widetilde{\mathfrak{X}}'(J)$ son image inverse dans
$\widetilde{\mathfrak{X}}'$. Soit $\mathcal{G}_{\mathfrak{X}'}(I,J)$ ou
simplement $\mathcal{G}'$ le sous-groupe de $\mathcal{G}_{\mathfrak{X}'}$ qui
invarie $J$. \struck{D. Voici les structures} $\mathcal{G}'$ opère sur $J$, sur
$\widetilde{\mathfrak{X}}'(J)$ et sur
\[
  (24)\qquad \widetilde{\mathfrak{X}}'(J) \longrightarrow J
\]
\note{« $\mathcal{G}'$ » est écrit au-dessus de la ligne biffée}

\page{8}
Enfin, on a une famille de scindages \marginal{interligne : « partiels »}
\[
  (25)\qquad \Psi = (\Psi_j)_{j \in J} : \bigl(\widetilde{\mathfrak{X}}(j)
  \longrightarrow \mathrm{Hom}(\mathcal{G}_j, \Sigma')\bigr)
\]
\note{au-dessus de $\Psi_j$ un petit mot biffé ; l'indice de
$\Sigma'$ est biffé}
(où cette fois $\Sigma'$ désigne l'image inverse
\marginal{interligne : « et $\mathcal{G}_j$ le stab.\ de $j$ dans
$\mathcal{G}'$ »} de $\mathcal{G}' \subset \mathcal{G}_{\mathfrak{X}'}$ dans
$\Sigma_{\mathfrak{X}'}$), compatible avec les opérations de
$\Sigma'$ de part et d'autre, \struck{s'inscrivant dans une famille de scindages} \struck{\ill{}}
\struck{mod $\pi_1$} \struck{(26)} \struck{$\Psi^{\mathrm{ext}} =
(\Psi_j^{\mathrm{ext}})_{j \in J}$} \struck{\ill{}}
\struck{$(\widetilde{\mathfrak{X}}'(J))$} \struck{$\to
\mathrm{Hom}(\mathcal{G}', \Sigma')/\mathrm{Int}(\pi_1, \Sigma')$}
\struck{\uncertain{équivariante}} \struck{compatible avec op.\ de $\Sigma'$}
\note{ce passage est barré de deux longs traits horizontaux et de traits
obliques}

\textcircled{6} \textbf{Récapitulons.} On a donc \struck{$\mathfrak{X}$ schéma
de t.f.\ sur $\overline{\mathbb{Q}}$, un diviseur}
\[
  \mathrm{I}\quad \left\{
  \begin{array}{l}
    \mathfrak{X} \text{ de t.f.\ sur } \overline{\mathbb{Q}} \\
    D = \sum_{i \in I} D_i \text{ diviseur effectif sur } \mathfrak{X}
      \text{ à comp.\ simples, } \mathfrak{X} \text{ lisse en les pts} \\
    \qquad \text{maximaux de } \operatorname{supp} D \\
    J \subset \mathfrak{X}'(\overline{\mathbb{Q}}), \text{ où }
      \mathfrak{X}' = \mathfrak{X} - D
  \end{array}\right.
\]
On pose
\[
  \mathrm{II}\quad \left\{
  \begin{array}{l}
    \mathcal{G}' = \text{groupe des autom.\ du schéma absolu } \mathfrak{X}
      \text{ qui invarie } D \text{ et } J \\
    G' = \text{s-gp de } \mathcal{G}' \text{ formé des }
      \overline{\mathbb{Q}}\text{-autom.} \\
    T(\overline{\mathbb{Q}}) = \varprojlim_n \mu_n(\overline{\mathbb{Q}}) \\
    \Gamma = \mathrm{Gal}(\overline{\mathbb{Q}}/\mathbb{Q})
  \end{array}\right.
\]
\note{« absolu » est ajouté au-dessus de la ligne ; sous le signe $\subset$ de
la dernière ligne de I, un petit mot souligné ne se lit pas}

\page{9}
\note{p.~5 de l'auteur. La numérotation des formules recommence ici à (1)}
d'où
\[
  (1)\qquad 1 \longrightarrow G' \longrightarrow \mathcal{G}' \longrightarrow \Gamma
\]
et un germe de relèvement

\begin{tikzcd}
(2)\quad \mathcal{G}' \arrow[r] & \Gamma \\
 & \Gamma' \arrow[ul] \arrow[u, hook]
\end{tikzcd}

\noindent($\Gamma' \subset \Gamma$ s-gp ouvert).
\marginal{en marge gauche, au crayon, en oblique, près de (2) : « pour mémoire :
relation avec \uncertain{variétés} \ill{} \ill{} \ill{}\ldots »}

\textbf{Choisissons} (cf.\ normalisation plus bas)

\textcircled{III} \textbf{un rev.\ universel} $\widetilde{\mathfrak{X}}'$
de $\mathfrak{X}'$, d'où \struck{posons}
\[
  \pi_1(\mathfrak{X}') = \mathrm{Aut}(\widetilde{\mathfrak{X}}'/\mathfrak{X}')
\]
On trouve alors
\[
  (3)\qquad \text{str.\ d'extension}\quad
  1 \to \pi_1(\mathfrak{X}') \longrightarrow \Sigma' \longrightarrow
  \mathcal{G}' \to 1
\]
\[
  (4)\qquad \Sigma'\text{-ensembles}\quad \left\{
  \begin{array}{l}
    \widetilde{\mathfrak{X}}'(I) \\
    \widetilde{\mathfrak{X}}'(J)
  \end{array}\right.
\]
\[
  (5)\qquad \text{Opérations de } \Sigma' \text{ sur } I, J \text{ via }
  \mathcal{G}', \text{ et } \Sigma'\text{-applications}\quad \left\{
  \begin{array}{l}
    \widetilde{\mathfrak{X}}'(I) \to I \\
    \widetilde{\mathfrak{X}}'(J) \to J \simeq \widetilde{\mathfrak{X}}'(J)/\pi_1
  \end{array}\right.
\]
\marginal{en marge gauche, au crayon, en oblique, près de (3) et (4) : « \ill{}
\ill{} $\mathfrak{X}(\overline{\mathbb{Q}})$ \ill{} \ill{}\ldots »}

\page{10}
\note{dans (6), un mot biffé précède $\widetilde{\mathfrak{X}}'(I)$}
\[
  (6)\qquad \text{application}\quad \Phi :
  \widetilde{\mathfrak{X}}'(I) \longrightarrow
  \mathrm{Hom}(T(\overline{\mathbb{Q}}), \pi_1(\mathfrak{X}'))
\]
compatible avec op.\ de $\Sigma'$, opérant sur le but de droite grâce à
son op.\ sur $T(\overline{\mathbb{Q}})$ via ($\mathcal{G}'$ via $\Gamma$) et sur
$\pi_1(\mathfrak{X}')$ [via autom.\ int.\ dans (3)]
\marginal{en marge gauche, au crayon, en oblique : « si $\dim \mathfrak{X} = 1$
sinon il faut revoir avec plus de soin la situation\ldots »}
\[
  (7)\qquad \text{Syst.\ d'application}\quad
  \Psi = (\Psi_j)_{j \in J} : \bigl(\widetilde{\mathfrak{X}}'(j)
  \longrightarrow \mathrm{Splitt}(\Sigma'_j / \mathcal{G}'_j)\bigr)_{j \in J}
\]
\note{au-dessus de $\Psi_j$ un petit signe biffé ; après la parenthèse un terme
$\mathrm{Int}(\pi_1, \Sigma')$ est biffé}
compatible avec opérations de $\Sigma'_j$ (opérant sur le but de droite
grâce aux autom.\ intérieurs
\[
  ({}^{g}u)(\sigma) = g\, u(g_1^{-1} \sigma g_1)\, g^{-1}
\]
où $\sigma \in \mathcal{G}'_j$, $u \in \mathrm{Splitt}(\Sigma'_j /
\mathcal{G}'_j)$, $g \in \Sigma'_j$, $g_1 \in \mathcal{G}'$ l'image de $g$
dans $\mathcal{G}'$).
\struck{(8)} \struck{$\Psi^{\mathrm{ext}} = (\Psi_j^{\mathrm{ext}})$ :}
\struck{$(\widetilde{\mathfrak{X}}'(J)$} \struck{$\to
\mathrm{Splitt}(\Sigma'_j / \mathcal{G}'_j)/\ldots)_{j \in J}$}
\note{la ligne (8) est biffée ; lecture incertaine dans son détail}

(9) germe d'homomorphisme

\begin{tikzcd}
\Sigma' \arrow[r] & \mathcal{G}' \arrow[r] & \Gamma \\
 & & \Gamma' \arrow[ull] \arrow[u, hook]
\end{tikzcd}

\noindent($\Gamma'$ s-groupe ouvert)
\marginal{en marge gauche, encadrée d'un trait, une note en oblique : « où
$\mathcal{G}'_j$ est le stabilisateur de $j$ dans $\mathcal{G}'$, $\Sigma'_j$
son image inverse dans $\Sigma'$ --- (8) \ill{} “\uncertain{relevant}” (7)
\ill{} $\mathcal{G}'$ \ill{} et \ill{} \ill{} $\Sigma'$\ldots »}

\page{11}
\note{p.~6 de l'auteur}
Programme des \struck{paradigmes de théorie des} \marginal{interligne :
« gruppentheoretische »} « Gleichnis » (paradigmes en termes de théorie des
groupes) de \struck{la théorie géométrique algébrique absolue (en car.\ 0)}
\marginal{interligne : « l'arithmétique »} \uncertain{l'arithmétique} :
\note{« Gleichnis » est en allemand, entre guillemets, dans l'original ; le mot
qui le précède, ajouté au-dessus de la ligne, est également allemand et se lit
mal}

a) Donner une description « \uncertain{gruppentheoretisch} » complète de
$\Gamma$

b) Donner une description complète de tous les éléments de structure énumérés
\struck{\ill{}} ci-dessus, dans des cas typiques (suivants pour
$(\mathfrak{X}, D, J)$), dont nous indiquons quelques-uns plus bas.
\marginal{en marge gauche, au crayon, en oblique : « On tiendra \ill{} de
\uncertain{prendre} $\mathfrak{X}$ projectif, \ill{} pour simplifier\ldots »}

\underline{NB} Comme on a les éléments de structure énumérés pour
$(\mathfrak{X}, D, J)$, on les a aussi \uncertain{automatiquement} pour
\marginal{interligne : « étale »} tout rev.\ fini de $\mathfrak{X}'$\ldots\ Il
faudrait expliciter ce point avec soin, et expliciter également dans les
\uncertain{équivalences} \uncertain{explicites} les \uncertain{fondements}
\uncertain{« bouclage de tresses »}. On rajoute \marginal{interligne :
« certains »} des $D_i$ (\uncertain{qu'on avait enlevés}) ou
« \uncertain{fixage} de \uncertain{tresses} » \marginal{sous la ligne :
« si $\dim \mathfrak{X} = 1$ »} (\uncertain{on enlève} certains des
$\xi \in J$).
\note{les guillemets de « \uncertain{bouclage de tresses} » et de
« \uncertain{fixage de tresses} » sont de lui ; les deux lectures sont
incertaines}

\page{12}
\underline{RemB} Soit $\Gamma'$ un s-groupe ouvert de $\Gamma$, correspondant à
une s-extension finie $K/\mathbb{Q}$ de $\overline{\mathbb{Q}}$, et considérons
une restriction de la situation $(\mathfrak{X}, D, J)$ à $K$, donnant donc par
un scindage partiel (cf.\ \textcircled{2})
\[
  \Gamma' \longrightarrow \mathcal{G}'
\]
d'où une extension
\[
  (9)\qquad 1 \to \pi_1(\mathfrak{X}') \longrightarrow \Sigma'_{\Gamma'}
  \longrightarrow \Gamma' \to 1
\]
Le groupe $\Gamma'$ opère donc sur $(\mathfrak{X}, D, J)$, donc sur
\marginal{interligne : « $\mathfrak{X}'$, donc sur »} la catégorie des rev.\
\marginal{interligne : « étales »} finis de $\mathfrak{X}'$, \struck{catégorie
qui} qui s'identifie à la cat.\ $\mathrm{Ens.f.}(\pi_1(\mathfrak{X}'))$. Pour
tout $\gamma \in \Gamma'$, soit $\mathcal{B}_{\gamma}$ le
$\pi_1(\mathfrak{X}')$-bitorseur image inverse de $\gamma$ dans
$\Sigma'_{\Gamma'}$, il \struck{définit} \marginal{interligne :
« explicite »} l'opération de $\gamma$ sur $\mathrm{Revét}(\mathfrak{X}')
\simeq \mathrm{Ensf}(\pi_1)$ comme
\[
  (10)\qquad E \longmapsto \mathcal{B}_{\gamma} \wedge^{\pi_1} E
\]
Si $\mathcal{Y}'$ et ${}^{\gamma}\mathcal{Y}'$ se correspondent, on a des isom.\
\struck{\ill{}} \marginal{interligne : « \uncertain{canoniques} »}
\[
  (11)\qquad \left\{
  \begin{array}{l}
    \mathcal{Y}'(I) \simeq ({}^{\gamma}\mathcal{Y}')(I) \\
    \mathcal{Y}'(J) = ({}^{\gamma}\mathcal{Y}')(J).
  \end{array}\right.
\]
\note{« RemB » est souligné ; lecture douteuse. La première ligne de (11) porte
$\simeq$, la seconde $=$, comme sur la page}

\page{13}
\note{p.~7 de l'auteur}
En fait, si $\mathcal{Y}' \longleftrightarrow E$, on a
\[
  \begin{array}{l}
    \mathcal{Y}'(I) \simeq \widetilde{\mathfrak{X}}'(I) \wedge^{\pi_1} E \\
    \mathcal{Y}'(J) \simeq \widetilde{\mathfrak{X}}'(J) \wedge^{\pi_1} E
  \end{array}
\]
et \struck{les appl.} isomorphismes (11) reviennent à des isom.
\[
  \left\{
  \begin{array}{l}
    \widetilde{\mathfrak{X}}'(I) \wedge^{\pi_1} E \simeq
      (\widetilde{\mathfrak{X}}'(I) \wedge^{\pi_1} \mathcal{B}_{\gamma}) \wedge^{\pi_1} E \\
    \widetilde{\mathfrak{X}}'(J) \wedge^{\pi_1} E \simeq
      (\widetilde{\mathfrak{X}}'(J) \wedge^{\pi_1} (\mathcal{B}_{\gamma})) \wedge^{\pi_1} E
  \end{array}\right.
\]
qui proviennent à \uncertain{droite} d'isom.\ de $\pi_1$-ens.\ à droite
\[
  \left\{
  \begin{array}{l}
    \widetilde{\mathfrak{X}}'(I) \simeq \widetilde{\mathfrak{X}}'(I) \wedge^{\pi_1} \mathcal{B}_{\gamma} \\
    \widetilde{\mathfrak{X}}'(J) \simeq \widetilde{\mathfrak{X}}'(J) \wedge^{\pi_1} \mathcal{B}_{\gamma}
  \end{array}\right.
\]
lesquels reviennent à la donnée d'homs
\[
  \left\{
  \begin{array}{l}
    \mathcal{B}_{\gamma} \longrightarrow \mathrm{Aut}(\widetilde{\mathfrak{X}}'(I))
      = \mathrm{Bij}(\widetilde{\mathfrak{X}}'(I), \widetilde{\mathfrak{X}}'(I)) \\
    \quad \text{et itou pour } \widetilde{\mathfrak{X}}'(J))
  \end{array}\right.
\]
compatibles avec les op.\ de $\pi_1$ à g.\ et à dr. On a donc là des
applications, pour $\gamma \in \Gamma'$ variable, équivalant \struck{\ill{}} à
la donnée d'\struck{une} applications
\[
  \Sigma'_{\Gamma'} = \bigcup_{\gamma \in \Gamma'} \mathcal{B}_{\gamma}
  \longrightarrow \mathrm{Aut}_{\mathrm{ens}}(\widetilde{\mathfrak{X}}'(I)),
  \qquad \longrightarrow \mathrm{Aut}_{\mathrm{ens}}(\widetilde{\mathfrak{X}}'(J))
\]
et celles qu'on cherche sont justement l'opération \uncertain{déduite} par
celles déjà envisagées de $\Sigma'$ sur $\widetilde{\mathfrak{X}}'(I)$,
$\widetilde{\mathfrak{X}}'(J)$\ldots
\note{sous « Aut », sur la troisième accolade, un mot surchargé ; l'indice
« ens » des deux derniers $\mathrm{Aut}$ est lu d'après le trait souligné}

\page{14}
\textcircled{7} Réconstruction\note{sic} des \uncertain{groupoïdes} fondamentaux
$\Pi_1(\mathfrak{X}'; J)$ de $\mathfrak{X}'$, et des opérations de
$\mathcal{G}'$ dessus, en termes de l'extension $\Sigma'$ de $\mathcal{G}'$
par $\pi_1$, et de l'opération de $\Sigma'$ sur
$\widetilde{\mathfrak{X}}'(J)$ et $\widetilde{\mathfrak{X}}'(J) \to J$ [et des
applications (7) (8)
\[
  \begin{array}{l}
    \Psi^{\mathrm{ext}} : \widetilde{\mathfrak{X}}'(J) \longrightarrow
      \mathrm{Splitt}(\Sigma'/\mathcal{G}')/\pi_1 \\
    \Psi_j : (\widetilde{\mathfrak{X}}'(j) \longrightarrow
      \mathrm{Splitt}(\Sigma'_j/\mathcal{G}'_j))
  \end{array}
\]
].
\marginal{en marge gauche, en oblique, séparée du texte par un trait : « inutiles
\uncertain{dans ces} \ill{}, \uncertain{elles résultent des précédentes} »}

\uncertain{Ob} $\Pi_1 = J$

\struck{$\mathfrak{X}'_j$, $j \in J$}
\underline{NB} Si $\widetilde{\xi} \in \widetilde{\mathfrak{X}}'(J)$,
$\widetilde{\mathfrak{X}}'[j]$ (rev.\ universel de $\mathfrak{X}'$ en $j$)
s'identifie à $\widetilde{\mathfrak{X}}'$ (puisque $\widetilde{\xi}$ est un
pt au-dessus de $j$) \ill{} des $\widetilde{\mathfrak{X}}'$ au dessus de $j$
\[
  \widetilde{\mathfrak{X}}' \xrightarrow{\ \psi_{\widetilde{\xi}}\ }
  \widetilde{\mathfrak{X}}'[j]
\]
Si $\gamma \in \pi_1$, on aura
\[
  \psi_{\gamma \cdot \widetilde{\xi}} = \psi_{\widetilde{\xi}} \circ
  \gamma_{\widetilde{\mathfrak{X}}'}
\]
donc on peut dire, indépendamment des choix \uncertain{d'un pt} base
$\widetilde{\xi}$
\[
  \widetilde{\mathfrak{X}}'[j] \xleftarrow{\ \sim\ }
  \underbrace{\widetilde{\mathfrak{X}}'(j)}_{\text{torseur à dr.\ sous } \pi_1}
  \wedge^{\pi_1} \widetilde{\mathfrak{X}}'
\]

\page{15}
\note{p.~8 de l'auteur}
Ceci noté, on a donc, si $j, j' \in J = \mathrm{Ob}\,\Pi_1$
\[
  \begin{aligned}
    \mathrm{Isom}_{\Pi_1}(j, j') &\overset{\mathrm{df}}{=}
      \mathrm{Isom}(\widetilde{\mathfrak{X}}'[j], \widetilde{\mathfrak{X}}'[j']) \\
    &\simeq \mathrm{Isom}_{\pi_1\text{-tors.\ à droite}}
      (\widetilde{\mathfrak{X}}'(j), \widetilde{\mathfrak{X}}'(j'))
  \end{aligned}
\]
la composition des chemins étant la composition des hom.\ de
$\pi_1$-torseurs à droite. \struck{\ill{}} On a donc reconstruit
\marginal{interligne : « le groupoïde »} $\Pi_1$ à l'aide de
$\widetilde{\mathfrak{X}}'(J)$ et de l'opération libre (à droite) de $\pi_1$ sur
$\widetilde{\mathfrak{X}}'(J)$. \struck{[Les opérations} Comment reconstituer
l'opération de $\mathcal{G}'$ sur $\Pi_1$ ? L'opération sur
$\mathrm{Ob}\,\Pi_1 = J$ n'est autre que l'opération donnée de $\mathcal{G}'$
sur $J$ (factorisant celle de $\Sigma'$). Soit $g \in \mathcal{G}'$,
explicitons
\[
  \begin{array}{ccc}
    \mathrm{Isom}_{\Pi_1}(j, j') & \xrightarrow{\ g_{\Pi_1}\ } &
      \mathrm{Isom}(gj, gj') \\
    \| & & \| \\
    \mathrm{Isom}_{\pi_1}(\widetilde{\mathfrak{X}}'(j), \widetilde{\mathfrak{X}}'(j'))
      & \overset{?}{\simeq} &
      \mathrm{Isom}_{\pi_1}(\widetilde{\mathfrak{X}}'(gj), \widetilde{\mathfrak{X}}'(gj'))
  \end{array}
\]
Pour ceci, considérons le \marginal{interligne : « $\pi_1$- »} bitorseur
\[
  \mathcal{B}_g \subset \Sigma'\quad \text{image inverse de } g \text{ par }
  \Sigma' \to \mathcal{G}'
\]
On note par l'\uncertain{hyp.}\ [grâce à l'op.\ de $\Sigma'$ sur
$\widetilde{\mathfrak{X}}'(I)$]
\[
  \begin{array}{l}
    \widetilde{\mathfrak{X}}'(gj) \simeq \widetilde{\mathfrak{X}}'(j) \wedge^{\pi_1} \mathcal{B}_g \\
    \widetilde{\mathfrak{X}}'(gj') \simeq \widetilde{\mathfrak{X}}'(j') \wedge^{\pi_1} \mathcal{B}_g
  \end{array}
\]
\note{dans le crochet, il écrit $\widetilde{\mathfrak{X}}'(I)$, là où le
contexte appellerait $\widetilde{\mathfrak{X}}'(J)$ ; transcrit tel quel}

\page{16}
d'où aussitôt la flèche envisagée.

Il faudrait encore expliciter comment (7) (8) résultent de la donnée de
l'opération de $\Sigma'$ sur $\widetilde{\mathfrak{X}}'(J)$, induisant une
opération \underline{libre} de $\pi_1 \subset \Sigma'$ sur
$\widetilde{\mathfrak{X}}'(J)$ (ayant $J$ comme quotient)\ldots

\underline{Reformulation}

Considérons les systèmes
\[
  \left\{
  \begin{array}{ll}
    1 \to \pi \longrightarrow \Sigma \longrightarrow G \to 1 &
      \text{ext.\ de groupes} \\
    \mathfrak{Z}\ \ \Sigma\text{-ens.}, \neq \emptyset,\ \text{où } \pi \subset
      \Sigma \text{ opère librement}
  \end{array}\right.
\]
\note{« $\neq \emptyset$ » est ajouté au-dessus de la ligne. Le groupe noté ici
$G$ est écrit avec la même lettre ronde que $\mathcal{G}'$ plus haut ; il est
rendu $G$, comme au second lot}
Ils forment une \underline{catégorie} $C$.

Considérons les systèmes $(\Pi, \widetilde{\Pi}, p, G, g)$ : $\Pi$ est un
groupoïde \marginal{interligne : « \uncertain{non vide} »} connexe,
$p : \widetilde{\Pi} \to \Pi$ un rev.\ universel de $\Pi$, et $G$ un groupe
opérant \struck{\ill{}} sur $\Pi$ par automorphismes par
$g : G \to \mathrm{Aut}(\Pi)$. Ils forment une catégorie $C'$. On va définir
des équivalences quasi-inverses l'une de l'autre
\[
  C \mathrel{\substack{\xrightarrow{\ \alpha\ } \\ \xleftarrow[\ \beta\ ]{}}} C'
\]

\page{17}
\note{p.~9 de l'auteur}
\underline{Définition de $\alpha$}. On va définir
\[
  \alpha(\Sigma, G, \pi, \mathfrak{Z}\ldots) = (\Pi, \widetilde{\Pi}, p, G, \sim)
\]
\[
  \left\{
  \begin{array}{l}
    \mathrm{Ob}\,\Pi = \mathfrak{Z}/\pi \quad (\text{soit } J) \\
    \mathrm{Isom}(j, j') = \mathrm{Isom}_{\pi}(\mathfrak{Z}_j, \mathfrak{Z}_{j'}) \\
    \text{composition des flèches évidente}
  \end{array}\right.
  \qquad \text{d'où groupoïde } \Pi
\]
\struck{$\widetilde{\Pi}_j = \mathfrak{Z}_j$} $\mathrm{Ob}\,\widetilde{\Pi} =
\mathfrak{Z}$
\[
  \begin{array}{l}
    \mathrm{Isom}(z, z') = \{\mathrm{pt}\} \quad \forall z, z' \in \mathfrak{Z}
      \quad (\widetilde{\Pi} \text{ cat.\ discrète\ldots}) \\
    \qquad\downarrow \\
    \mathrm{Isom}(j, j') \simeq \mathrm{Isom}(\mathfrak{Z}_j, \mathfrak{Z}_{j'})
      \quad (\text{si } j = \bar z,\ j' = \bar z')
  \end{array}
\]
\marginal{à droite, séparé par un trait vertical : « $\widetilde{\Pi}$ cat.\
\underline{discrète} sur $\Pi$ définie par \struck{\ill{}} \uncertain{un}
foncteur $F : \Pi^{\circ} \to \mathrm{Ens}$, $F(j) = \mathfrak{Z}_j$, \ldots,
$\widetilde{\Pi} \simeq \Pi_{/F}$ »}
\note{sous la parenthèse « (si $j = \bar z$, $j' = \bar z'$) », une ligne biffée
ne se lit pas}
$\downarrow$ transforme l'unique pt de $\mathrm{Isom}(z, z')$ en l'unique
\marginal{interligne : « $\pi$- »} isom.\ $\mathfrak{Z}_j \to \mathfrak{Z}_{j'}$
transformant $z$ en $z'$\ldots

Pour l'instant, on n'a pas utilisé $\Sigma$ et l'op.\ de $\Sigma$ sur
$\mathfrak{Z}$. On va l'utiliser maintenant pour faire opérer
$G = \Sigma/\pi$ sur $\Pi$, et $\Sigma$ sur $\widetilde{\Pi}$, de façon
compatible. Comme $\Sigma$ opère sur $\mathfrak{Z} = \mathrm{Ob}(\widetilde{\Pi})$,
il opère \struck{\ill{}} sur $\widetilde{\Pi}$, cat.\ discrète associée,
\struck{\ill{}} et \struck{\ill{}} $\pi$ opère par \uncertain{$\Pi$-automorphismes},
d'où l'opération de $G$ sur $\Pi$ par passage au quotient.

\page{18}
\underline{Définition de $\beta$}. On définit
\[
  \beta(\Pi, \widetilde{\Pi}, p, G, \sim) = (\Sigma, G, \pi, \mathfrak{Z})
  \quad \text{ainsi :}
\]
$\Sigma$ = groupe des couples $(g, \bar g)$, avec $g \in G$,
$\bar g \in \mathrm{Aut}(\widetilde{\Pi})$ compatible avec $g$
\[
  \pi = \mathrm{Aut}(\widetilde{\Pi}/\Pi)
\]
$\Sigma \to G$ par $(g, \bar g) \mapsto g$, d'où
\[
  1 \to \pi \longrightarrow \Sigma \longrightarrow G \to 1
\]
$\mathfrak{Z} = \mathrm{Ob}\,\widetilde{\Pi}$, avec opération évidente de
$\Sigma$ sur $\mathfrak{Z}$\ldots

Variantes profinies\ldots

\textcircled{8} Considérons, avec notations de \textcircled{7},
($1 \to \pi \to \Sigma \to G \to 1$ et $\mathfrak{Z}$ $\Sigma$-ens.\
$\neq \emptyset$ sur lequel $\pi$ opère librement)
$\longleftrightarrow (\Pi, \widetilde{\Pi}, p, G, \sim)$. On va interpréter la
donnée d'un triple $(\mathcal{T}, T, \Phi)$, où $\mathcal{T}$ est
\marginal{interligne : « $\chi$ : »} un $\Sigma$-ens., $T$
\struck{\ill{}} un groupe \uncertain{\ill{}} ($\simeq \widehat{\mathbb{Z}}$ ?),
$\chi : G \to \mathrm{Aut}(T)$ ($\simeq \widehat{\mathbb{Z}}^{*}$
\uncertain{\ill{}}), et
\[
  \Phi : \mathcal{T} \longrightarrow \mathrm{Hom}_{\mathrm{gr}}(T, \pi)
\]
compatible avec op.\ de $\Sigma$ (opérant
\note{la lettre $\mathfrak{Z}$ (une sorte de « 3 » à queue) désigne l'ensemble
qui correspond, dans le récapitulatif de la page~8, à
$\widetilde{\mathfrak{X}}'(J)$ ; la lettre $\mathcal{T}$ (un T ronde barré),
celui qui correspond à $\widetilde{\mathfrak{X}}'(I)$ et à l'application $\Phi$
de (6), p.~10. Devant le premier $\mathcal{T}$ du triple, un petit signe en
exposant ne se lit pas}

\page{19}
\note{p.~10 de l'auteur}
sur $\mathcal{T}$ par l'op.\ donnée, et sur $\mathrm{Hom}_{\mathrm{gr}}(T, \pi)$ en
considérant $T$, $\pi$ comme des $\Sigma$-groupes grâce à $\chi$ et à
l'opération de $\Sigma$ sur $\pi$ par autom.\ intérieurs\ldots\ Pour
\struck{ceci, introduisons}
\note{à partir d'ici et jusqu'au haut de la page~20, le texte est traversé par
deux longs traits obliques parallèles, qui semblent le biffer ; il reste lisible
et est transcrit}
\[
  J = \mathcal{T}/\pi
\]
sur lequel \marginal{interligne : « et $=$ \ldots\ $G$ »} $\Sigma/\pi$
opère par passage au quotient, $\mathcal{T} \to J$ devenant donc un
$\Sigma$-morphisme. Considérons alors le groupoïde $\langle T, J \rangle$
défini par
\[
  \begin{array}{l}
    \mathrm{Ob}\,\langle T, J \rangle = J \\
    \mathrm{Isom}(j, j') = \left\{
      \begin{array}{ll} \emptyset & \text{si } j \neq j' \\
        T & \text{si } j = j' \end{array}\right.
  \end{array}
\]
composition des flèches étant celle de $T$. On a une opération de $G$ sur
$\langle T, J \rangle$ via ses opérations sur $T$ \marginal{interligne :
« (par $\chi$) »} et sur $J$. On va définir un homomorphisme de groupoïdes
\[
  \langle T, J \rangle \longrightarrow \Pi^{\wedge}
\]
où $\Pi^{\wedge}$ est le groupoïde \marginal{interligne : « 0-connexe »} des
rev.\ universels de $\Pi$, identifiable (via $\widetilde{\Pi}$, qui en est un
$\in \Pi^{\wedge}$) au groupoïde des

\page{20}
$\pi$-torseurs. Pour $j \in J$, le $\pi$-torseur correspondant $P_j$ s'envoie
dans $\mathcal{T}_j$ = fibre de $j \in \mathcal{T}/\pi$ dans $\mathcal{T}$,
\struck{puis} la fibre de $P_j$ en $t \in \mathcal{T}_j$ étant le torseur sous
le stabilisateur $\pi_t$ de $t$ dans $\pi$ suivant : on considère
\[
  \Phi(t) : T \longrightarrow \pi
\]
dont $\pi_t$ centralise
\note{fin du passage traversé par les deux traits obliques}

Considérons la restriction de l'opération de $\Sigma$ sur $\mathcal{T}$ à
$\pi$, d'où un groupoïde $\langle \mathcal{T}, \pi \rangle$ dont l'ens.\ des
objets est $\mathcal{T}$, l'ens.\ des $\mathrm{Isom}(t, t')$ étant
$\pi_{t \to t'} = \{\lambda \in \pi \mid \lambda t = t'\}$.
\struck{La composition} \ill{}.
\[
  \pi_0 \langle \mathcal{T}, \pi \rangle = \mathcal{T}/\pi \overset{\mathrm{df}}{=} J
\]
\[
  \pi_1(\mathcal{T}, t) = \pi_t \quad (\text{stabilisateur de } t \text{ dans } \pi)
\]
et on trouve que $\pi_t$ \struck{centralise} $\Phi(t)(T) \subset \pi$ : \ill{},
\uncertain{supposons} \underline{$T$ commutatif}
et \struck{$\pi_t = \Phi_t(T)$}
\[
  \Phi_t : T \xrightarrow{\ \sim\ } \pi_t
\]
\struck{du} d'autres termes, on suppose donnée, en plus de $\mathcal{T}$
$\Sigma$-ens.\ \uncertain{de $\pi$-ens.},
\note{la phrase se poursuit au début du second lot (p.~21)}

\end{document}
